% SPDX-License-Identifier: Apache-2.0
\section{An Array of Children}
\label{sec:x-lanes}

A unit of units whose children are a count of one unit keeps them as
one field, \code{lanes: Units<Lane, N>}, and joins them with
\code{join\_all} over \code{iter\_mut().enumerate()}, giving child
\code{i} its ends by index (issue 635). \code{chans::<T, C, N>()} makes
the \code{N} channels between them as two sets of ends, and
\code{Ends::from} makes a set of an array port; \code{take(e)} hands
an end out once, where \code{e} is \code{i} or a neighbour's index.

\code{Lanes<N>} is a ring. Each lane sends its input to the lane on
its right, keeps what the lane on its left sent, and puts out its
input plus that. Lane \code{i} sends on \code{ring\_i} and receives on
\code{ring\_\{(i + N - 1) \% N\}}. The ring is of channels: a word on a
channel is seen the step after it is sent, whichever lane runs first,
where a wire read in the step it is driven would depend on the order.
The run asserts that every lane puts out its own input plus its left
neighbour's from two cycles before. The netlist of \code{Lanes<3>},
three instances and three channels, is checked against the run under
nvc and Verilator, as Section~\ref{sec:x-checked} describes.

\begin{figure}[htbp]
\centering
\input{lanes_timing.tex}
\caption{Three lanes in a ring: each keeps what its left neighbour
sent, and puts out its own input plus that.}
\label{fig:x-lanes}
\end{figure}

\lstinputlisting[caption={\code{ex\_lanes.rs}. Run with \code{bazel run //lib/examples:ex\_lanes}.},label={lst:x-lanes}]{lib/examples/ex_lanes.rs}

\lstinputlisting[style=ascii,caption={What \code{ex\_lanes} printed when the build ran it: the outputs, the instances, then the lowering.},label={lst:x-lanes-out}]{out_lanes.txt}
